%<*driver>
\def\nameofplainTeX{plain}
\ifx\fmtname\nameofplainTeX
\input l3docstrip.tex
\keepsilent
\askforoverwritefalse
\preamble

========================================================================
 semcat - generalized semantic-category markup with optional QR codes.

 (C) 2026 Matthias Werner
 E-mail: matthias.werner@informatik.tu-chemnitz.de

 Released under the LaTeX Project Public License v1.3c or later
 See http://www.latex-project.org/lppl.txt
========================================================================

\endpreamble
\postamble

 This work is "maintained" (as per LPPL maintenance status) by
 Matthias Werner.
\endpostamble
\generate{\file{semcat.sty}{\from{semcat.dtx}{package}}}
\expandafter\endbatchfile
\fi

\iffalse
%</driver>
%<*package>
\NeedsTeXFormat{LaTeX2e}[2022-06-01]
\ProvidesExplPackage{semcat}{2026-08-31}{v0.3.0}
  {Generalized semantic-category markup with optional QR annotations}
\ExplSyntaxOff
%</package>
%<*driver>
\fi

\documentclass[add-bib=false]{osgdoc}

\OnlyDescription

\usepackage{enumitem}
\usepackage{microtype}
\usepackage[
languages={de,en},
map={de=german,en=english},
targetlang={job=2},
load babel
]{langselect}
\newpackagename{\semcat}{semcat}
\makeindex
\setcnltx{
  package  =  semcat,
  version  = \UseVersionOf{semcat.sty},
  date     = \UseDateOf{semcat.sty},
  title    = {\ldeen{Das \semcat\ Paket}{The \semcat\ Package}},
  info     = {\ldeen{Semantische Kategorienmarkierung mit optionalen QR-Codes}{
      Semantic category markup with optional QR codes}},
  authors  = {Matthias Werner[matthias.werner@informatik.tu-chemnitz.de]},
  abstract={
    \ldeen*{
      Das Paket ordnet Textabschnitte benannten, farblich unterschiedenen
      Kategorien zu -- etwa "`Pflichtstoff"', "`Vertiefung"' oder
      "`Hintergrund"' -- und stellt dafür Inline-Markierung, eine über
      mehrere Absätze und Seitenumbrüche hinweg tragfähige Block-Markierung
      (Box oder durchgehender Randbalken), eine schwebende Erklärbox sowie
      QR-Code-Annotationen bereit. Kategorien sind frei definierbar; ist
      @1 geladen, können Farbe und Darstellung aus dessen Paletten und
      Dekorationen stammen. Sind @2 aktiv, folgen Sichtbarkeit und
      Darstellung den Modi; bei aktivem @3 erzeugen die Markierungen
      eigene PDF-Tag-Strukturen.

      Das Paket ist Teil des \pkg*{osglecture}-Bundles, hängt aber -- bis
      auf diese beiden optionalen Integrationen -- von keinem anderen
      Paket des Bundles ab und ist eigenständig nutzbar.
    }{
      This package assigns passages of text to named, color-distinguished
      categories -- for instance "required", "supplementary", or
      "background" -- and provides inline marking, a block-level marking
      that survives multiple paragraphs and page breaks (a box or a
      continuous margin bar), a floating explanation box, and QR-code
      annotations for them. Categories are freely definable; when @1 is
      loaded, color and rendering may come from its palettes and
      decorations. When @2 is active, visibility and rendering follow
      the active modes; when @3 is active, the markup produces its own
      PDF tag structures.

      This package is part of the \pkg*{osglecture} bundle, but -- apart
      from those two optional integrations -- depends on none of the
      bundle's other packages and can be used on its own.
    }{\pkg{osgstyler}}{\pkg{osglecture-modes}}{\cs{DocumentMetadata}\Marg{tagging=on}}
  },
  url      =https://github.com/werner-matthias/osglecture,
  build-title
}
\begin{document}
% --------------------------------------------------------------------------
% Documentation
% --------------------------------------------------------------------------
\section{\ldeen{Nutzung}{Usage}}

\subsection{\ldeen{Kategorien}{Categories}}

\ldeen{
  Vier Kategorien -- \code{A} (Standardkategorie), \code{B}, \code{C} und
  \code{D} -- sind vordefiniert. Eigene Kategorien legt \cs{SemCatDefine}
  an:
}{
  Four categories -- \code{A} (the default), \code{B}, \code{C} and
  \code{D} -- are predefined. \cs{SemCatDefine} declares further
  categories:
}

\begin{verbatim}
\SemCatDefine{E}[label={Historical background}, color=teal]
\SemCatDefine*{F}[label={Core material}, slot=primary]
\end{verbatim}

\ldeen*{
  Die Sternform macht die Kategorie zur neuen Standardkategorie. Ohne
  @1-Schlüssel wird die Farbe reihum aus der eigens dafür vorgesehenen
  Palette \code{semcat-categories} vergeben, sofern \pkg{osgstyler}
  geladen ist, sonst aus einer eingebauten Farbfolge. Diese Palette ist
  bewusst von den generischen Theme-Slots getrennt: Ein Themenwechsel
  im Dokument soll die Kategoriefarben nicht mitverändern. Ihre
  Vorgabewerte stammen aus ColorBrewer~"`Dark2"'; wer andere
  Kategoriefarben möchte, deklariert \code{semcat-categories} nach dem
  Laden von semcat einfach erneut mit \cs{DeclareOsgColorPalette}. Der
  @2-Schlüssel wählt gezielt einen Slot dieser Palette -- wirkt aber
  nur, solange \pkg{osgstyler} geladen ist; ohne das Paket greift
  wieder die eingebaute Rotation.
}{
  The star form makes the category the new default. Without a @1 key,
  the color is assigned round-robin from the dedicated
  \code{semcat-categories} palette while \pkg{osgstyler} is loaded,
  otherwise from a built-in color sequence. This palette is
  deliberately kept separate from the generic theme slots: switching
  the document's theme should not also change category colors. Its
  defaults come from ColorBrewer "Dark2"; anyone wanting different
  category colors simply declares \code{semcat-categories} again
  after loading semcat, via \cs{DeclareOsgColorPalette}. The @2 key
  picks a slot of this palette explicitly -- but only while
  \pkg{osgstyler} is loaded; without it, the built-in rotation applies
  instead.
}{\code{color}}{\code{slot}}

\ldeen*{
  Zusätzlich zu Farbe/Slot kann eine Kategorie festlegen, \emph{wie}
  ihre Blockmarkierung (siehe unten) dargestellt wird: als Box
  (Vorgabe) oder als durchgehender Randbalken. Ist \pkg{osgstyler}
  geladen, kann stattdessen eine dort deklarierte Dekoration
  referenziert werden -- deren \code{margin}-Komponente liefert dann
  den Randbalken, ihre \code{enclosure}-Komponente die Box, jeweils mit
  eigenen Farben/Maßen:
}{
  In addition to color/slot, a category can specify \emph{how} its
  block-level marking (see below) renders: as a box (the
  default) or as a continuous margin bar. While \pkg{osgstyler} is
  loaded, a decoration declared there can be referenced instead -- its
  \code{margin} component then supplies the bar, its \code{enclosure}
  component the box, each with its own colors/measurements:
}

\begin{verbatim}
\SemCatDefine{G}[label={Aside}, style=bar]
\DeclareOsgDecoration{marked-caution}{
  enclosure = { border-color=accent, fill-color=secondary }
}
\SemCatDefine{H}[label={Caution}, decoration=marked-caution]
\end{verbatim}

\ldeen*{
  Ein @1-Schlüssel überschreibt Stil oder Dekoration abhängig vom
  aktiven osglecture-Modus -- zusätzlich zum unqualifizierten
  Standardwert, der als Rückfallebene dient, wenn kein aktiver Modus
  (oder gar kein @2) einen Eintrag hat:
}{
  A @1 key overrides style or decoration depending on the active
  osglecture mode -- on top of the unqualified default, which serves
  as the fallback when no active mode (or no @2 at all) has an entry:
}{\code{variants}}{\pkg{osglecture-modes}}

\begin{verbatim}
\SemCatDefine{I}[
  label={Aside}, style=box,
  variants={ slides = {style=bar} }
]
\end{verbatim}

\ldeen*{
  Das Konzept ist bisher stark auf Farben ausgerichtet -- beim Druck
  eines Lehrbuchs stehen aber oft keine oder wenige Farben zur
  Verfügung. Für diesen Fall kann ein Symbol an die Stelle der
  Kategoriefarbe treten, am Anfang des Randbalkens bzw.\ im Kopf der
  Box: Der @1-Schlüssel benennt ein registriertes Symbol,
  \cs{SemCatDeclareSymbol} registriert weitere (die eingebauten
  stammen aus \pkg{pifont}, das -- anders als Unicode-Symbole -- in
  praktisch jeder Installation verfügbar ist). Ist eine
  osgstyler-Dekoration mit einer \code{affix}-Komponente aufgelöst,
  hat deren \code{content}- bzw.\ \code{symbol}-Schlüssel Vorrang vor
  dem einfachen @1-Schlüssel der Kategorie; \code{position},
  \code{font} und \code{gap} der \code{affix}-Komponente werden derzeit
  nicht ausgewertet.
}{
  So far the concept is strongly color-oriented -- but printing a
  textbook often leaves few or no colors available. For that case, a
  symbol can take the category color's place, at the start of the
  margin bar or in the box header: the @1 key names a registered
  symbol, and \cs{SemCatDeclareSymbol} registers further ones (the
  built-in ones come from \pkg{pifont}, which -- unlike Unicode
  symbols -- is available in virtually every installation). When a
  resolved osgstyler decoration has an \code{affix} component, its
  \code{content}/\code{symbol} key takes precedence over the
  category's plain @1 key; the \code{affix} component's
  \code{position}, \code{font} and \code{gap} are not evaluated yet.
}{\code{symbol}}

\begin{verbatim}
\SemCatDeclareSymbol{pin}{\ding{43}}
\SemCatDefine{J}[label={Caution}, symbol=warning]
\end{verbatim}

\subsection{\ldeen{Markierung}{Marking}}

\ldeen{
  \cs{SemCat} markiert Text inline mit der Kategoriefarbe -- unabhängig
  vom für die Kategorie gewählten \code{style}, da Inline-Text nicht
  sinnvoll als Box oder Balken erscheinen kann:
}{
  \cs{SemCat} marks text inline with the category color -- independent
  of the category's chosen \code{style}, since inline text cannot
  sensibly appear as a box or a bar:
}

\begin{verbatim}
Ordinary text with a \SemCat{B}{marked phrase} in the middle.
\end{verbatim}

\ldeen{
  Für ganze Abschnitte -- auch über mehrere Absätze und Seitenumbrüche
  hinweg -- markiert die Umgebung \code{SemCatMark} den Inhalt gemäß
  dem für die Kategorie (und den aktiven Modus) aufgelösten Stil.
  Ein Randbalken wird dabei nicht wie beim früheren, changebar-basierten
  Vorgänger über die Output-Routine gezeichnet, sondern zeilenweise über
  einen \LuaLaTeX-Callback -- dadurch übersteht er Seitenumbrüche ohne
  Sonderbehandlung:
}{
  For whole passages -- even across multiple paragraphs and page
  breaks -- the \code{SemCatMark} environment marks the content
  according to the style resolved for the category (and the active
  mode). A margin bar is not drawn via the output routine the way the
  former, changebar-based predecessor did, but line by line via a
  \LuaLaTeX{} callback -- so it survives page breaks without special
  handling:
}

\begin{verbatim}
\begin{SemCatMark}[category=B]
  Supplementary material, with a footnote\footnote{like this one},
  spanning as many paragraphs as needed.
\end{SemCatMark}
\end{verbatim}

\subsection{\ldeen{Erklärbox}{Explanation box}}

\ldeen{
  \code{SemCatExplain} ist -- anders als \cs{SemCat} und
  \code{SemCatMark}, die immer \emph{im} Lesefluss bleiben -- eine
  echte, nummerierte und beschriftete Gleitumgebung für Inhalt
  \emph{außerhalb} des Leseflusses: überspringbare Zusatzerklärungen.
  Sie ist strukturell stets eine Box, unabhängig vom sonst für die
  Kategorie gewählten Stil, und nutzt die Kategoriefarbe nur für ihre
  optische Wiedererkennbarkeit:
}{
  Unlike \cs{SemCat} and \code{SemCatMark}, which always stay
  \emph{within} the reading flow, \code{SemCatExplain} is a genuine,
  numbered and captioned floating environment for content \emph{outside}
  the reading flow: skippable supplementary explanations. It is
  structurally always a box, regardless of the style otherwise chosen
  for the category, and uses the category color only for visual
  identity:
}

\begin{verbatim}
\begin{SemCatExplain}[category=B]{Why this matters}
  Supplementary background that a reader could skip.
\end{SemCatExplain}
\end{verbatim}

\ldeen{
  Die Umgebung basiert auf \pkg{float} und registriert den Gleittyp
  \code{semcatexplain} (Vorgabeplatzierung \code{tbp}, Verzeichnisendung
  \code{loe}); \cs{listof}\Marg{semcatexplain}\Marg{...} erzeugt ein
  Verzeichnis aller Erklärboxen.
}{
  The environment is built on \pkg{float} and registers the
  \code{semcatexplain} float type (default placement \code{tbp}, list
  extension \code{loe}); \cs{listof}\Marg{semcatexplain}\Marg{...}
  produces a list of all explanation boxes.
}

\subsection{\ldeen{QR-Code-Annotationen}{QR-code annotations}}

\ldeen{
  \cs{SemCatQR} setzt einen QR-Code mit Beschriftung, standardmäßig als
  Randnotiz, mit dem Stern inline:
}{
  \cs{SemCatQR} places a QR code with a caption, by default as a margin
  note, inline with the star form:
}

\begin{verbatim}
\SemCatQR[B]{https://example.org}{Further reading}
\SemCatQR*[B]{https://example.org}{Further reading}
\end{verbatim}

\ldeen{
  \cs{SemCatQRRedirect} baut die Ziel-URL stattdessen aus einer zuvor
  mit \cs{SemCatSetRedirectBase} gesetzten Basis-URL und einer
  Dokument-/Eintrags-ID zusammen; unter \LuaLaTeX{} werden beide Teile
  robust URL-kodiert:
}{
  \cs{SemCatQRRedirect} instead builds the target URL from a base URL set
  via \cs{SemCatSetRedirectBase} and a document/entry id; under
  \LuaLaTeX{} both parts are robustly URL-encoded:
}

\begin{verbatim}
\SemCatSetRedirectBase{https://example.org/redirect.php}
\SemCatQRRedirect[B]{42}{Further reading}
\end{verbatim}

\subsection{\ldeen{Modus- und Tagging-Integration}{Mode and tagging integration}}

\ldeen*{
  Alle Markierungsbefehle akzeptieren eine führende
  Winkelklammer-Spezifikation, die -- ist @1 geladen -- an
  \cs{IfLectureModeTF} weitergereicht wird; ohne @1 wird sie ignoriert
  und die Markierung erscheint unbedingt:
}{
  Every marking command accepts a leading angle-bracket specification
  that -- while @1 is loaded -- is forwarded to \cs{IfLectureModeTF};
  without @1 it is ignored and the markup always appears:
}{\pkg{osglecture-modes}}

\begin{verbatim}
\SemCat<article>{B}{printed-only remark}
\end{verbatim}

\ldeen*{
  \code{SemCatMark} und \code{SemCatExplain} erhalten den Modus
  stattdessen über den @1-Schlüssel, da sie als \pkg*{environ}-Umgebungen
  keine Winkelklammer-Syntax entgegennehmen können; ohne @1-Schlüssel
  erscheint die Markierung unbedingt, unabhängig von aktiven Modi:
}{
  \code{SemCatMark} and \code{SemCatExplain} instead take the mode via
  the @1 key, since as \pkg*{environ} environments they cannot accept
  angle-bracket syntax; without a @1 key, the markup always appears,
  regardless of active modes:
}{\code{mode}}

\begin{verbatim}
\begin{SemCatMark}[mode=slides,category=B]
  Slides-only supplementary passage.
\end{SemCatMark}
\end{verbatim}

\ldeen{
  Ist ein Dokument mit \cs{DocumentMetadata}\Marg{tagging=on} getaggt,
  registriert jede Kategorie eigene Strukturnamen (\code{semcat/<Key>}
  als \code{Span}, \code{semcat/<Key>/note} als \code{Note}), und
  Inline-Markierung sowie \code{SemCatMark} werden entsprechend getaggt
  -- beide bleiben \emph{im} Lesefluss. \code{SemCatExplain} nutzt
  dagegen die einzige, kategorieunabhängige Struktur
  \code{semcat/explain} mit der Rolle \code{Aside}, da sie bewusst
  \emph{außerhalb} des Leseflusses steht. QR-Codes selbst bleiben
  ungetaggt, da \pkg{qrcode} sie über eine eigene \pkg*{tikz}-Grafik
  zeichnet, die sich nicht sauber in eine manuell geöffnete
  Tag-Struktur einfügt; die Beschriftung daneben wird normal als
  Fließtext getaggt.
}{
  When a document is tagged via \cs{DocumentMetadata}\Marg{tagging=on},
  every category registers its own structure names (\code{semcat/<key>}
  as \code{Span}, \code{semcat/<key>/note} as \code{Note}), and inline
  marking and \code{SemCatMark} are tagged accordingly -- both stay
  \emph{within} the reading flow. \code{SemCatExplain} instead uses the
  single, category-independent structure \code{semcat/explain} with
  role \code{Aside}, since it deliberately sits \emph{outside} the
  reading flow. QR codes themselves stay untagged, since \pkg{qrcode}
  draws them via its own \pkg*{tikz} graphic, which does not nest
  cleanly inside a manually opened tag structure; the caption next to
  it is tagged normally as ordinary flow text.
}

\section{Implementation}
\MaybeStop{\ldeen*{Um die Implementationsbeschreibung zu erhalten, kommentieren Sie bitte \cs{OnlyDescription} in @1 aus.}{To view the implementation description, please uncomment @1 in \code{osgdoc.dtx}.}{\file{semcat.dtx}}}

\subsection{\emph{semcat.sty}}
\AutoDocInput{semcat.dtx}{package}
\subsection{semcat.lua}
\AutoDocInput[lua]{semcat.lua}{lua}
\end{document}

%</driver>

%<*package>
\RequirePackage[dvipsnames]{xcolor}
\RequirePackage{luacolor}
\RequirePackage{qrcode}
\RequirePackage{marginnote}
\RequirePackage{footnotehyper}
\RequirePackage{environ}
\RequirePackage{pifont}
\RequirePackage{float}
\RequirePackage{tcolorbox}
\tcbuselibrary{skins,breakable}
\newfloat{semcatexplain}{tbp}{loe}
\floatname{semcatexplain}{Explanation}
\ExplSyntaxOn

% ---------------------------------------------------------------------
% Category registry
% ---------------------------------------------------------------------
\seq_new:N \g__semcat_categories_seq
\prop_new:N \g__semcat_category_label_prop
\tl_new:N \g__semcat_default_category_tl
\tl_new:N \l__semcat_tmp_tl
\tl_new:N \l__semcat_tmp_color_tl

\int_new:N \g__semcat_fallback_slot_int
\clist_new:N \c__semcat_fallback_colors_clist
\clist_set:Nn \c__semcat_fallback_colors_clist
  { NavyBlue, BurntOrange, OliveGreen, RedViolet, Turquoise, Plum }

% \ldeen{Statt die reihum vergebenen Kategoriefarben aus dem
% Aktiv-Theme zu ziehen (\cs{OsgColorName}, an das jeweils
% \emph{aktive} \cs{UseOsgColorPalette} gebunden), deklariert semcat
% eine eigene, feste Palette \code{semcat-categories}: Ein
% Themenwechsel im Dokument soll die Kategoriefarben nicht
% mitverändern, und die Slots eines gewöhnlichen Themes (\code{primary}
% etc.) sind auf Harmonie \emph{innerhalb} eines Themes ausgelegt, nicht
% auf gegenseitige Unterscheidbarkeit als Kategoriensatz. Die
% Vorgabewerte stammen aus ColorBrewer~"`Dark2"' (eine verbreitete,
% für gute Unterscheidbarkeit ausgelegte qualitative Palette). Wer
% andere Kategoriefarben möchte, deklariert \code{semcat-categories}
% nach dem Laden von semcat einfach erneut mit
% \cs{DeclareOsgColorPalette} -- eine erneute Deklaration desselben
% Namens ist in osgstyler ausdrücklich vorgesehen. \code{text} und
% \code{background} bleiben absichtlich unbelegt: sie sind für
% Fließtext-/Seitenfarbe reserviert und ungeeignet als
% Kategorieakzent, und diese Palette ist ohnehin nicht dazu gedacht,
% selbst per \cs{UseOsgColorPalette} aktiviert zu werden.}{Rather than
% drawing the round-robin category colors from the active theme
% (\cs{OsgColorName}, tied to whichever \cs{UseOsgColorPalette} is
% \emph{active}), semcat declares its own, fixed palette
% \code{semcat-categories}: switching the document's theme should not
% also change category colors, and an ordinary theme's slots
% (\code{primary} etc.) are tuned for harmony \emph{within} one theme,
% not for mutual distinguishability as a category set. The defaults
% come from ColorBrewer "Dark2" (a widely used qualitative palette
% chosen for distinguishability). Anyone wanting different category
% colors simply declares \code{semcat-categories} again after loading
% semcat, via \cs{DeclareOsgColorPalette} -- redeclaring the same name
% is explicitly supported by osgstyler. \code{text} and
% \code{background} are deliberately left unset: they are reserved for
% body-text/page color and unsuited as a category accent, and this
% palette is not meant to ever be activated itself via
% \cs{UseOsgColorPalette} anyway.}
\clist_new:N \c__semcat_palette_slots_clist
\clist_set:Nn \c__semcat_palette_slots_clist
  { primary, secondary, accent, structure, alert, example }
\tl_const:Nn \c__semcat_palette_name_tl { semcat-categories }

\cs_new_protected:Npn \__semcat_declare_default_palette:
  {
    \exp_args:No \DeclareOsgColorPalette { \c__semcat_palette_name_tl }
      {
        primary   = { model = HTML, value = 1B9E77 },
        secondary = { model = HTML, value = D95F02 },
        accent    = { model = HTML, value = 7570B3 },
        structure = { model = HTML, value = E7298A },
        alert     = { model = HTML, value = 66A61E },
        example   = { model = HTML, value = E6AB02 }
      }
  }
\IfPackageLoadedTF { osgstyler }
  { \__semcat_declare_default_palette: }
  {
    \AddToHook { package/osgstyler/after } [ semcat/palette ]
      { \__semcat_declare_default_palette: }
  }

% \ldeen{Definiert (bzw. redefiniert, falls schon vorhanden) unter
% einem fest aus #1/#2 gebildeten Namen eine xcolor-Farbe mit dem
% aktuellen Wert des Slots #2 der Palette #1 -- absichtlich bei jedem
% Aufruf frisch aus \cs{OsgColorPaletteValue} gelesen, statt einmalig
% zwischengespeichert, damit eine spätere Neudeklaration von #1
% (z.\,B. durch den Nutzer) auch rückwirkend wirkt, solange sie vor
% dem jeweiligen \cs{SemCatDefine}-Aufruf erfolgt.}{Defines (or
% redefines, if already existing) an xcolor color, under a name built
% deterministically from #1/#2, holding the current value of palette
% #1's slot #2 -- deliberately read fresh from \cs{OsgColorPaletteValue}
% on every call rather than cached once, so that a later redeclaration
% of #1 (e.g. by the user) still takes effect, as long as it happens
% before the relevant \cs{SemCatDefine} call.}
\cs_new_protected:Npn \__semcat_palette_color:nnN #1#2#3
  {
    \tl_set:Nx \l__semcat_tmp_tl { \OsgColorPaletteValue {#1} {#2} {model} }
    \tl_set:Nx \l__semcat_tmp_color_tl { \OsgColorPaletteValue {#1} {#2} {value} }
    \tl_set:Nx #3 { semcat-palettecolor-#2 }
    \exp_args:Nooo \definecolor #3 { \l__semcat_tmp_tl } { \l__semcat_tmp_color_tl }
  }

\msg_new:nnn { semcat } { unknown-category }
  { Unknown~semantic~category~'#1'. }
\msg_new:nnn { semcat } { category-redefined }
  { Semantic~category~'#1'~is~already~defined;~redefining~it. }
\msg_new:nnn { semcat } { redirect-base-missing }
  { No~QR~redirect~base~URL~is~configured;~call~SemCatSetRedirectBase~first. }
\msg_new:nnn { semcat } { luatex-required-for-redirect }
  { SemCatQRRedirect~needs~LuaLaTeX~to~encode~the~redirect~target~robustly;~falling~back~to~the~unencoded~value~'#1'. }
\msg_new:nnn { semcat } { unknown-symbol }
  { Unknown~semantic~symbol~'#1'. }

% \ldeen{Ein kleines, eigenständiges Symbolregister: Wenn keine Farbe
% verfügbar ist (typischerweise beim Druck eines Lehrbuchs), kann ein
% Symbol an die Stelle der Kategoriefarbe treten -- am Anfang des
% Randbalkens bzw.\ im Kopf der Box. osgstylers Dekorationsmodell sieht
% dafür bereits die \code{affix}-Komponente vor
% (\code{content}/\code{symbol}), liefert aber selbst keinen
% "`Symbol-Backend"'-Mechanismus, der \code{symbol}-Referenzen in
% tatsächliche Glyphen übersetzt (osgstylers eigene Dokumentation weist
% ausdrücklich darauf hin, dass aktuell kein Backend dafür existiert).
% semcat übernimmt diese Rolle selbst, mit \pkg{pifont} als Backend --
% universell verfügbar, ohne Font-Abdeckungsrisiko, im Gegensatz zu
% Unicode-Symbolen, die in Fließtextschriften oft fehlen.}{A small,
% self-contained symbol registry: when no color is available
% (typically when printing a textbook), a symbol can take the
% category color's place -- at the start of the margin bar, or in the
% box header. osgstyler's decoration model already provides for this
% via the \code{affix} component (\code{content}/\code{symbol}), but
% does not itself ship a "symbol backend" mechanism that resolves
% \code{symbol} references into actual glyphs (osgstyler's own
% documentation explicitly notes that no such backend currently
% exists). semcat takes on that role itself, using \pkg{pifont} as the
% backend -- universally available, with no font-coverage risk, unlike
% Unicode symbols, which are often missing from body-text fonts.}
\prop_new:N \g__semcat_symbol_prop
\NewDocumentCommand \SemCatDeclareSymbol { m m }
  { \prop_gput:Nnn \g__semcat_symbol_prop {#1} {#2} }
\NewExpandableDocumentCommand \SemCatSymbol { m }
  { \prop_item:Nn \g__semcat_symbol_prop {#1} }
\cs_new_protected:Npn \__semcat_symbol_lookup:nN #1#2
  {
    \prop_if_in:NnTF \g__semcat_symbol_prop {#1}
      { \tl_set:No #2 { \prop_item:Nn \g__semcat_symbol_prop {#1} } }
      {
        \msg_warning:nnn { semcat } { unknown-symbol } {#1}
        \tl_clear:N #2
      }
  }
\SemCatDeclareSymbol{check}{\ding{51}}
\SemCatDeclareSymbol{cross}{\ding{55}}
\SemCatDeclareSymbol{warning}{\ding{55}}
\SemCatDeclareSymbol{caution}{\ding{55}}
\SemCatDeclareSymbol{star}{\ding{72}}
\SemCatDeclareSymbol{important}{\ding{72}}
\SemCatDeclareSymbol{pencil}{\ding{46}}
\SemCatDeclareSymbol{note}{\ding{46}}
\SemCatDeclareSymbol{hand}{\ding{37}}
\SemCatDeclareSymbol{attention}{\ding{37}}
\SemCatDeclareSymbol{square}{\ding{110}}
\SemCatDeclareSymbol{arrow}{\ding{242}}
\SemCatDeclareSymbol{reference}{\ding{242}}

% \ldeen{Prüft einen bereits aufgelösten Kategorieschlüssel in \#1 und
% ersetzt ihn -- nach Meldung -- durch die Standardkategorie, falls er
% nie registriert wurde. So erzeugt ein falscher Schlüssel genau eine
% diagnostizierbare Fehlermeldung statt einer Kaskade von
% Folgefehlern ("`undefined color"', "`unknown tag"') aus jeder
% nachfolgenden Abfrage, die einen gültigen Schlüssel voraussetzt.}
% {Validates an already-resolved category key held in \#1 and, after
% reporting it, substitutes the default category if it was never
% registered. This way one bad key produces a single diagnosable error
% rather than a cascade of "undefined color"/"unknown tag" errors from
% every later lookup that assumed a valid key.}
\cs_new_protected:Npn \__semcat_validate_category:N #1
  {
    \exp_args:NNo \seq_if_in:NnF \g__semcat_categories_seq #1
      {
        \msg_error:nnx { semcat } { unknown-category } {#1}
        \tl_set_eq:NN #1 \g__semcat_default_category_tl
      }
  }

\cs_new:Npn \__semcat_color_name:n #1 { semcat-#1 }
\NewExpandableDocumentCommand \SemCatColor { m } { \__semcat_color_name:n {#1} }
\NewExpandableDocumentCommand \SemCatLabel { m }
  { \prop_item:Nn \g__semcat_category_label_prop {#1} }
\NewExpandableDocumentCommand \SemCatDefaultCategory { }
  { \g__semcat_default_category_tl }

% \ldeen{Löst ein optionales Kategorieargument (leer oder ein Makro)
% gegen die registrierte Standardkategorie zu einem einfachen, voll
% expandierten Schlüssel auf, damit jede folgende Abfrage mit einer
% konkreten Zeichenkette arbeitet statt mit Tokens, die noch expandiert
% werden müssten.}{Resolves an optional category argument (possibly
% empty, possibly a macro) against the registered default into a plain,
% fully expanded key, so every later lookup works against a concrete
% string rather than tokens that would still need expanding.}
\cs_new_protected:Npn \__semcat_resolve_category:nN #1#2
  {
    \tl_set:Nx #2 {#1}
    \tl_if_empty:NT #2 { \tl_set_eq:NN #2 \g__semcat_default_category_tl }
  }

% \ldeen{Alles Folgende nimmt einen bereits aufgelösten, einfachen
% Kategorieschlüssel (eine Kontrollsequenz wie \cs{l_tmpa_tl} mit
% Inhalt "`A"') entgegen und erzeugt daraus über einen expliziten
% \cs{tl_set:Nx}-Schritt ein weiteres, voll expandiertes Ergebnis:
% keiner dieser Helfer verlässt sich darauf, dass ein Verbraucher
% (xcolor, tcolorbox, die Tagging-Schlüssel des Kerns) eine
% Variablenreferenz selbst expandiert.}{Everything below takes an
% already-resolved plain category key (a control sequence such as
% \cs{l_tmpa_tl} holding e.g. "A") and produces a further fully
% expanded result via one explicit \cs{tl_set:Nx} step, so these
% helpers never rely on a downstream consumer (xcolor, tcolorbox, the
% kernel's tagging keys) to expand a variable reference itself.}
\cs_new_protected:Npn \__semcat_resolve_color:nN #1#2
  { \tl_set:Nx #2 { \__semcat_color_name:n {#1} } }

\cs_new:Npn \__semcat_structure_name:n #1 { semcat/#1 }
\cs_new:Npn \__semcat_note_structure_name:n #1 { semcat/#1/note }

\cs_new_protected:Npn \__semcat_tag_begin:n #1
  { \cs_if_exist:NT \tagstructbegin { \tagstructbegin{ tag=#1 } } }
\cs_new_protected:Npn \__semcat_tag_end:
  { \cs_if_exist:NT \tagstructend { \tagstructend } }
% \ldeen{Eine einzelne, kategorieunabhängige Struktur für die
% Erklärbox: Anders als Inline-Markierung und Randmarkierung, die immer
% \emph{im} Lesefluss bleiben (Rolle \code{Span}/\code{Note}), steht
% die Erklärbox bewusst \emph{außerhalb} davon -- als eigenständiges,
% überspringbares Float. Die PDF-Rolle \code{Aside} bildet genau das
% ab.}{A single, category-independent structure for the explanation
% box: unlike inline marking and the margin marker, which always stay
% \emph{within} the reading flow (role \code{Span}/\code{Note}), the
% explanation box deliberately sits \emph{outside} it -- a standalone,
% skippable float. The PDF role \code{Aside} captures exactly that.}
\cs_if_exist:NT \NewStructureName
  {
    \NewStructureName{ semcat/explain }
    \AssignStructureRole{ semcat/explain } { Aside }
  }
% \ldeen*{Baut einen voll aufgelösten Aufruf
% \cs{UseStructureName}\Marg{semcat/<Schlüssel>} (einen konkreten,
% bereits expandierten Strukturnamen, nie eine Variablenreferenz) und
% übergibt ihn unausgewertet an \cs{__semcat_tag_begin:n}, analog dazu,
% wie \file{ansiterm.dtx} \cs{UseStructureName}-Aufrufe an den
% @1-Schlüssel von \cs{tagstructbegin} übergibt.}{Builds a
% fully-resolved \cs{UseStructureName}\Marg{semcat/<key>} call (a
% concrete, already-expanded structure name, never a variable
% reference) and hands it unevaluated to \cs{__semcat_tag_begin:n},
% mirroring how \file{ansiterm.dtx} feeds \cs{UseStructureName} calls to
% \cs{tagstructbegin}'s @1 key.}{\code{tag=}}
\cs_new_protected:Npn \__semcat_tag_begin_named:n #1
  { \__semcat_tag_begin:n { \UseStructureName{#1} } }
\cs_new_protected:Npn \__semcat_tag_begin_span:n #1
  {
    \tl_set:Nx \l__semcat_tmp_tl { \__semcat_structure_name:n {#1} }
    \exp_args:No \__semcat_tag_begin_named:n \l__semcat_tmp_tl
  }
\cs_new_protected:Npn \__semcat_tag_begin_note:n #1
  {
    \tl_set:Nx \l__semcat_tmp_tl { \__semcat_note_structure_name:n {#1} }
    \exp_args:No \__semcat_tag_begin_named:n \l__semcat_tmp_tl
  }

% \ldeen*{Ohne @1/@2 wird die Farbe reihum aus der Palette
% \code{semcat-categories} vergeben, wenn osgstyler geladen ist, sonst
% aus @3. Ein Slotname wird nur verwendet, wenn osgstyler tatsächlich
% geladen ist; andernfalls fällt die Zuordnung auf die eingebaute
% Rotation zurück.}{Without @1/@2, the color is assigned round-robin
% from the \code{semcat-categories} palette when osgstyler is loaded,
% otherwise from @3. A slot name is only honored while osgstyler is
% actually loaded; otherwise assignment falls back to the built-in
% rotation.}{color}{slot}{\cs{clist}\Arg{NavyBlue,...}}
\cs_new_protected:Npn \__semcat_next_color:N #1
  {
    \int_gincr:N \g__semcat_fallback_slot_int
    \IfPackageLoadedTF { osgstyler }
      {
        \exp_args:Nxx \__semcat_palette_color:nnN { \c__semcat_palette_name_tl }
          {
            \clist_item:Nn \c__semcat_palette_slots_clist
              {
                \int_mod:nn { \g__semcat_fallback_slot_int - 1 }
                  { \clist_count:N \c__semcat_palette_slots_clist } + 1
              }
          }
          #1
      }
      {
        \tl_set:Nx #1
          {
            \clist_item:Nn \c__semcat_fallback_colors_clist
              {
                \int_mod:nn { \g__semcat_fallback_slot_int - 1 }
                  { \clist_count:N \c__semcat_fallback_colors_clist } + 1
              }
          }
      }
  }

\cs_if_exist:NT \NewStructureName
  {
    \cs_new_protected:Npn \__semcat_register_tags:n #1
      {
        \NewStructureName{ semcat/#1 }
        \AssignStructureRole{ semcat/#1 } { Span }
        \NewStructureName{ semcat/#1/note }
        \AssignStructureRole{ semcat/#1/note } { Note }
      }
  }
\cs_if_exist:NF \NewStructureName
  { \cs_new_protected:Npn \__semcat_register_tags:n #1 { } }

\keys_define:nn { semcat / category }
  {
    color      .tl_set:N  = \l__semcat_category_color_tl,
    color      .initial:n = { },
    slot       .tl_set:N  = \l__semcat_category_slot_tl,
    slot       .initial:n = { },
    label      .tl_set:N  = \l__semcat_category_label_tl,
    label      .initial:n = { },
    style      .tl_set:N  = \l__semcat_category_style_tl,
    style      .initial:n = { },
    decoration .tl_set:N  = \l__semcat_category_decoration_tl,
    decoration .initial:n = { },
    symbol     .tl_set:N  = \l__semcat_category_symbol_tl,
    symbol     .initial:n = { },
    variants   .code:n    = { \keys_set:nn { semcat / category / variants } {#1} },
  }

\NewDocumentCommand \SemCatDefine { s m O{} }
  {
    \seq_if_in:NnTF \g__semcat_categories_seq {#2}
      { \msg_warning:nnn { semcat } { category-redefined } {#2} }
      {
        \seq_gput_right:Nn \g__semcat_categories_seq {#2}
        \__semcat_register_tags:n {#2}
      }
    \tl_set:Nn \l__semcat_category_being_defined_tl {#2}
    \tl_clear:N \l__semcat_category_color_tl
    \tl_clear:N \l__semcat_category_slot_tl
    \tl_clear:N \l__semcat_category_label_tl
    \tl_clear:N \l__semcat_category_style_tl
    \tl_clear:N \l__semcat_category_decoration_tl
    \tl_clear:N \l__semcat_category_symbol_tl
    \keys_set:nn { semcat / category } {#3}
    \tl_if_empty:NF \l__semcat_category_style_tl
      { \prop_gput:Nee \g__semcat_style_variant_prop { #2 / default } \l__semcat_category_style_tl }
    \tl_if_empty:NF \l__semcat_category_decoration_tl
      { \prop_gput:Nee \g__semcat_decoration_variant_prop { #2 / default } \l__semcat_category_decoration_tl }
    \tl_if_empty:NF \l__semcat_category_symbol_tl
      { \prop_gput:Nee \g__semcat_symbol_variant_prop { #2 / default } \l__semcat_category_symbol_tl }
    \tl_if_empty:NTF \l__semcat_category_label_tl
      { \prop_gput:Nnn \g__semcat_category_label_prop {#2} {#2} }
      { \prop_gput:NnV \g__semcat_category_label_prop {#2} \l__semcat_category_label_tl }
    \tl_if_empty:NT \l__semcat_category_color_tl
      {
        \tl_if_empty:NTF \l__semcat_category_slot_tl
          { \__semcat_next_color:N \l__semcat_category_color_tl }
          {
            \IfPackageLoadedTF { osgstyler }
              { \tl_set:Nx \l__semcat_category_color_tl
                  { \OsgColorName { \l__semcat_category_slot_tl } } }
              { \__semcat_next_color:N \l__semcat_category_color_tl }
          }
      }
    \colorlet { \__semcat_color_name:n {#2} } { \l__semcat_category_color_tl }
    \IfBooleanT {#1} { \tl_gset:Nn \g__semcat_default_category_tl {#2} }
  }

% \ldeen{Berücksichtigt eine Modusspezifikation nur, solange
% osglecture-modes (oder ein kompatibles Backend) tatsächlich
% \cs{IfLectureModeTF} bereitstellt; sonst ist die Einschränkung
% wirkungslos, sodass semcat nie hart davon abhängt.}{Honors a
% mode-spec only while osglecture-modes (or a compatible backend)
% actually provides \cs{IfLectureModeTF}; otherwise the restriction is
% a no-op, so semcat never hard-depends on it.}
\cs_new_protected:Npn \__semcat_if_mode:nn #1#2
  {
    \tl_if_empty:nTF {#1}
      {#2}
      {
        \cs_if_exist:NTF \IfLectureModeTF
          { \IfLectureModeT {#1} {#2} }
          {#2}
      }
  }

% ---------------------------------------------------------------------
% Inline marking
% ---------------------------------------------------------------------
\NewDocumentCommand \SemCat { D<>{} m +m }
  {
    \__semcat_if_mode:nn {#1}
      {
        \group_begin:
        \__semcat_resolve_category:nN {#2} \l_tmpa_tl
        \__semcat_validate_category:N \l_tmpa_tl
        \exp_args:No \__semcat_tag_begin_span:n \l_tmpa_tl
        \exp_args:No \__semcat_resolve_color:nN \l_tmpa_tl \l__semcat_tmp_color_tl
        \color{\l__semcat_tmp_color_tl}#3%
        \__semcat_tag_end:
        \group_end:
      }
  }
\cs_if_exist:NT \MakeOverlayAwareCommand
  {
    \AddToHook { package/osglecture-modes/after } [ semcat/modes ]
      { \MakeOverlayAwareCommand [ arguments = { m +m } ] { \SemCat } }
  }

% ---------------------------------------------------------------------
% Presentation resolution: which style (box/bar) -- and, while osgstyler
% is loaded, which of its decorations -- renders a category's
% block-level marking. Resolved from the category AND the active
% lecture mode(s) together, since the same category may need a boxed
% callout on slides but a plain margin bar in the printed script.
% ---------------------------------------------------------------------
\prop_new:N \g__semcat_style_variant_prop
\prop_new:N \g__semcat_decoration_variant_prop
\prop_new:N \g__semcat_symbol_variant_prop
\tl_new:N \l__semcat_variant_key_tl
\tl_new:N \l__semcat_category_being_defined_tl
\tl_new:N \l__semcat_resolved_style_tl
\tl_new:N \l__semcat_resolved_decoration_tl
\tl_new:N \l__semcat_resolved_symbol_tl
\clist_new:N \l__semcat_mode_clist

\cs_new_protected:Npn \__semcat_variant_lookup:NnN #1#2#3
  {
    \tl_set:Nx \l__semcat_variant_key_tl {#2}
    \exp_args:NNo \prop_if_in:NnTF #1 \l__semcat_variant_key_tl
      { \tl_set:Nx #3 { \exp_args:NNo \prop_item:Nn #1 \l__semcat_variant_key_tl } }
      { \tl_clear:N #3 }
  }

% \ldeen{Ein \code{variants}-Schlüssel bei \cs{SemCatDefine} nimmt eine
% verschachtelte Schlüsselliste entgegen, deren Schlüssel Modusnamen
% sind und deren Werte wiederum @1/@2 setzen können -- diese
% Modusüberschreibungen landen in denselben Eigenschaftslisten wie der
% unqualifizierte Standardwert, nur unter dem jeweiligen Modusnamen
% statt "`default"'.}{A \code{variants} key on \cs{SemCatDefine} takes
% a nested key list whose keys are mode names and whose values may in
% turn set @1/@2 -- these mode overrides land in the same property
% lists as the unqualified default value, just under that mode's name
% instead of "default".}{\code{style}}{\code{decoration}}
\keys_define:nn { semcat / variant }
  {
    style      .tl_set:N  = \l__semcat_variant_style_tl,
    style      .initial:n = { },
    decoration .tl_set:N  = \l__semcat_variant_decoration_tl,
    decoration .initial:n = { },
    symbol     .tl_set:N  = \l__semcat_variant_symbol_tl,
    symbol     .initial:n = { },
  }
\cs_new_protected:Npn \__semcat_variant_declare:nn #1#2
  { % #1 = mode name, #2 = keyval content for that mode
    \tl_clear:N \l__semcat_variant_style_tl
    \tl_clear:N \l__semcat_variant_decoration_tl
    \tl_clear:N \l__semcat_variant_symbol_tl
    \keys_set:nn { semcat / variant } {#2}
    \tl_if_empty:NF \l__semcat_variant_style_tl
      { \prop_gput:Nee \g__semcat_style_variant_prop
          { \l__semcat_category_being_defined_tl / #1 } \l__semcat_variant_style_tl }
    \tl_if_empty:NF \l__semcat_variant_decoration_tl
      { \prop_gput:Nee \g__semcat_decoration_variant_prop
          { \l__semcat_category_being_defined_tl / #1 } \l__semcat_variant_decoration_tl }
    \tl_if_empty:NF \l__semcat_variant_symbol_tl
      { \prop_gput:Nee \g__semcat_symbol_variant_prop
          { \l__semcat_category_being_defined_tl / #1 } \l__semcat_variant_symbol_tl }
  }
\keys_define:nn { semcat / category / variants }
  { unknown .code:n = { \__semcat_variant_declare:nn { \l_keys_key_str } {#1} } }

% \ldeen*{Versucht für die Kategorie #1 der Reihe nach jeden aktiven
% Modus (bei geladenem osglecture-modes) und zuletzt "`default"', bis
% Stil, Dekoration und Symbol gefunden sind; ohne osglecture-modes
% bleibt nur "`default"' übrig, also genau der beim
% \cs{SemCatDefine}-Aufruf hinterlegte Wert.}{Tries, for category #1,
% each active mode in turn (while osglecture-modes is loaded) and
% finally "default", until style, decoration and symbol are found;
% without osglecture-modes only "default" remains, i.e. exactly the
% value recorded at \cs{SemCatDefine} time.}
\cs_new_protected:Npn \__semcat_resolve_presentation:nNNN #1#2#3#4
  {
    \clist_clear:N \l__semcat_mode_clist
    \cs_if_exist:NT \ActiveLectureModes
      { \clist_set:Nx \l__semcat_mode_clist { \ActiveLectureModes } }
    \clist_put_right:Nn \l__semcat_mode_clist { default }
    \tl_clear:N #2
    \tl_clear:N #3
    \tl_clear:N #4
    \clist_map_inline:Nn \l__semcat_mode_clist
      {
        \tl_if_empty:NT #3
          { \__semcat_variant_lookup:NnN \g__semcat_decoration_variant_prop { #1 / ##1 } #3 }
        \tl_if_empty:NT #2
          { \__semcat_variant_lookup:NnN \g__semcat_style_variant_prop { #1 / ##1 } #2 }
        \tl_if_empty:NT #4
          { \__semcat_variant_lookup:NnN \g__semcat_symbol_variant_prop { #1 / ##1 } #4 }
      }
  }

\seq_new:N \l__semcat_metric_seq
% \ldeen{Übersetzt einen Dekorationswert wie \code{rule/thin} über
% \cs{OsgMetric} in ein konkretes Maß; ein Wert ohne Schrägstrich gilt
% bereits als fertiges Maß.}{Translates a decoration value such as
% \code{rule/thin} via \cs{OsgMetric} into a concrete measurement; a
% value without a slash already counts as a finished measurement.}
\cs_new_protected:Npn \__semcat_metric_resolve:nN #1#2
  {
    \seq_set_split:Nnn \l__semcat_metric_seq { / } {#1}
    \int_compare:nNnTF { \seq_count:N \l__semcat_metric_seq } = { 2 }
      {
        \tl_set:Nx #2
          {
            \OsgMetric
              { \seq_item:Nn \l__semcat_metric_seq {1} }
              { \seq_item:Nn \l__semcat_metric_seq {2} }
          }
      }
      { \tl_set:Nn #2 {#1} }
  }

% ---------------------------------------------------------------------
% Margin bar rendering: a continuous, page-break-safe rule in the left
% margin, drawn per output line by a LuaTeX post_linebreak_filter
% callback (semcat.lua) rather than by hooking the output routine the
% way changebar does. \g__semcat_bar_box holds a plain \vrule, freshly
% colored per use; the callback clones it per line and only needs to
% know which box register to read, so no color/dimension state has to
% cross the Lua boundary beyond that one integer.
% ---------------------------------------------------------------------
\newbox\g__semcat_bar_box
\dim_new:N \l__semcat_bar_thickness_dim

% \ldeen{\pkg{luacolor} (oben angefordert) macht \cs{color} über
% Zeilen- und Seitenumbrüche hinweg sicher; ohne dieses Paket würde die
% hier gesetzte Farbe in Folgetext ausbluten, sobald die Box, aus der
% \g__semcat_bar_box geklont wird, mehrfach über das Dokument verteilt
% verwendet wird.}{\pkg{luacolor} (required above) makes \cs{color}
% safe across line and page breaks; without it, the color set here
% would bleed into later text once the box \g__semcat_bar_box is
% cloned from is reused repeatedly throughout the document.}
\cs_new_protected:Npn \__semcat_bar_setup:n #1
  { % #1 = resolved, plain xcolor color name
    \global\setbox\g__semcat_bar_box=\hbox{\color{#1}\vrule width \l__semcat_bar_thickness_dim}
    \directlua{ semcat.bar_begin(\the\g__semcat_bar_box) }
  }
\cs_new_protected:Npn \__semcat_bar_teardown:
  { \directlua{ semcat.bar_end() } }

\tl_new:N \l__semcat_enclosure_frame_tl
\tl_new:N \l__semcat_enclosure_back_tl
\tl_new:N \l__semcat_enclosure_rule_tl
\tl_new:N \l__semcat_enclosure_radius_tl

% ---------------------------------------------------------------------
% Affix (icon) resolution: an osgstyler decoration's \code{affix}
% component takes precedence when present; otherwise a category's
% plain \code{symbol} key (also mode-varied via \code{variants}, same
% as style/decoration) supplies the marker. Either way, the resolved
% marker is placed at the start of the margin bar or in the box
% header, so it can carry the category's identity when color alone is
% not distinguishable enough (or not available at all, as in a
% monochrome print run).
% ---------------------------------------------------------------------
\tl_new:N \l__semcat_affix_content_tl
\tl_new:N \l__semcat_affix_color_tl

% \ldeen{Liest Inhalt (\code{content} vor \code{symbol}, wie bei
% osgstylers eigener Affix-Verarbeitung) und Farbe aus der
% \code{affix}-Komponente der benannten Dekoration #1; ohne eigene
% Farbe gilt die Kategoriefarbe.}{Reads content (\code{content} before
% \code{symbol}, matching osgstyler's own affix handling) and color
% from the named decoration #1's \code{affix} component; without an
% explicit color, the category color applies.}
\cs_new_protected:Npn \__semcat_affix_from_decoration:n #1
  {
    \tl_set:Nx \l__semcat_tmp_tl { \OsgDecorationValue {#1} { affix } { content } }
    \tl_if_empty:NTF \l__semcat_tmp_tl
      {
        \tl_set:Nx \l__semcat_tmp_tl { \OsgDecorationValue {#1} { affix } { symbol } }
        \tl_if_empty:NF \l__semcat_tmp_tl
          { \exp_args:No \__semcat_symbol_lookup:nN \l__semcat_tmp_tl \l__semcat_affix_content_tl }
      }
      { \tl_set_eq:NN \l__semcat_affix_content_tl \l__semcat_tmp_tl }
    \tl_set:Nx \l__semcat_tmp_tl { \OsgDecorationValue {#1} { affix } { color } }
    \tl_if_empty:NF \l__semcat_tmp_tl
      { \tl_set:Nx \l__semcat_affix_color_tl { \OsgColorName { \l__semcat_tmp_tl } } }
  }

% \ldeen{Löst den Affix für die aktuelle Markierung auf: zuerst die
% \code{affix}-Komponente der aufgelösten Dekoration (falls vorhanden),
% sonst das aufgelöste, einfache \code{symbol} der Kategorie/Variante.
% Ohne beides bleibt der Affix leer, und weder Balken noch Box zeigen
% einen Marker.}{Resolves the affix for the current marking: first the
% resolved decoration's \code{affix} component (if present), else the
% category/variant's resolved plain \code{symbol}. Without either, the
% affix stays empty and neither the bar nor the box shows a marker.}
\cs_new_protected:Npn \__semcat_resolve_affix:
  {
    \tl_clear:N \l__semcat_affix_content_tl
    \tl_set_eq:NN \l__semcat_affix_color_tl \l__semcat_box_color_tl
    \tl_if_empty:NF \l__semcat_resolved_decoration_tl
      {
        \IfPackageLoadedTF { osgstyler }
          {
            \exp_args:No \OsgDecorationIfHasComponentTF \l__semcat_resolved_decoration_tl { affix }
              { \exp_args:No \__semcat_affix_from_decoration:n \l__semcat_resolved_decoration_tl }
              { }
          }
          { }
      }
    \tl_if_empty:NT \l__semcat_affix_content_tl
      {
        \tl_if_empty:NF \l__semcat_resolved_symbol_tl
          { \exp_args:No \__semcat_symbol_lookup:nN \l__semcat_resolved_symbol_tl \l__semcat_affix_content_tl }
      }
  }

\cs_new_protected:Npn \__semcat_maybe_affix_marginnote:
  {
    \tl_if_empty:NF \l__semcat_affix_content_tl
      { \marginnote{\textcolor{\l__semcat_affix_color_tl}{\l__semcat_affix_content_tl}}[0pt] }
  }

\tl_new:N \l__semcat_box_title_tl
\cs_new_protected:Npn \__semcat_compose_title:
  {
    \tl_if_empty:NTF \l__semcat_affix_content_tl
      { \tl_set:Nx \l__semcat_box_title_tl { \exp_args:No \SemCatLabel \l__semcat_box_category_tl } }
      {
        \tl_set:Nn \l__semcat_box_title_tl
          {
            \textcolor { \l__semcat_affix_color_tl } { \l__semcat_affix_content_tl } ~
            \exp_args:No \SemCatLabel \l__semcat_box_category_tl
          }
      }
  }

% \ldeen{Ein \cs{directlua}-Aufruf ist ein Whatsit-Knoten, der ausgeführt
% wird, sobald ihn der Eingabestrom erreicht -- nicht erst, wenn der
% umgebende Absatz tatsächlich in Zeilen umbrochen wird. Endet @1 nicht
% mit einer Leerzeile, bliebe der letzte Absatz beim Erreichen dieses
% Aufrufs noch offen und würde erst durch das abschließende \cs{end}
% umgebrochen -- zu diesem Zeitpunkt wäre die Balkenfarbe längst
% zurückgesetzt. Das erzwungene \cs{par} schließt den Absatz sofort,
% während der Balken noch aktiv ist.}{A \cs{directlua} call is a whatsit
% node that executes as soon as the input stream reaches it -- not
% when the surrounding paragraph actually gets broken into lines. If
% @1 does not end with a blank line, the last paragraph would still be
% open when this call is reached and would only be broken later by the
% closing \cs{end} -- by which time the bar color would already have
% been reset. The forced \cs{par} closes the paragraph immediately,
% while the bar is still active.}{\cs{BODY}}
\cs_new_protected:Npn \__semcat_render_bar_default:
  {
    \dim_set:Nn \l__semcat_bar_thickness_dim { 2pt }
    \exp_args:No \__semcat_bar_setup:n \l__semcat_box_color_tl
    \__semcat_maybe_affix_marginnote:
    \BODY
    \par
    \__semcat_bar_teardown:
  }

% \ldeen{Liest Farbe und Dicke aus der \code{margin}-Komponente der
% benannten osgstyler-Dekoration #1; \code{side}, \code{pattern} und
% \code{offset} werden derzeit nicht ausgewertet, der Balken erscheint
% immer links mit durchgezogener Linie.}{Reads color and thickness from
% the named osgstyler decoration #1's \code{margin} component;
% \code{side}, \code{pattern}, and \code{offset} are not evaluated yet,
% the bar always appears on the left with a solid line.}
\cs_new_protected:Npn \__semcat_render_bar_from_decoration:n #1
  {
    \tl_set:Nx \l__semcat_tmp_tl { \OsgDecorationValue {#1} { margin } { color } }
    \tl_if_empty:NTF \l__semcat_tmp_tl
      { \tl_set_eq:NN \l__semcat_tmp_color_tl \l__semcat_box_color_tl }
      { \tl_set:Nx \l__semcat_tmp_color_tl { \OsgColorName { \l__semcat_tmp_tl } } }
    \tl_set:Nx \l__semcat_tmp_tl { \OsgDecorationValue {#1} { margin } { thickness } }
    \tl_if_empty:NTF \l__semcat_tmp_tl
      { \dim_set:Nn \l__semcat_bar_thickness_dim { 2pt } }
      {
        \exp_args:No \__semcat_metric_resolve:nN \l__semcat_tmp_tl \l__semcat_tmp_tl
        \exp_args:No \dim_set:Nn \l__semcat_bar_thickness_dim { \l__semcat_tmp_tl }
      }
    \exp_args:No \__semcat_bar_setup:n \l__semcat_tmp_color_tl
    \__semcat_maybe_affix_marginnote:
    \BODY
    \par
    \__semcat_bar_teardown:
  }

\cs_new_protected:Npn \__semcat_render_box_default:
  {
    \__semcat_compose_title:
    \begin{tcolorbox}[
        enhanced, breakable,
        colback = \l__semcat_box_color_tl!12,
        colframe = \l__semcat_box_color_tl!70!black,
        fonttitle = \sffamily,
        title = {\l__semcat_box_title_tl},
        arc = 3pt, boxrule = .4pt,
        \l__semcat_box_options_tl
      ]
      \BODY
    \end{tcolorbox}
  }

% \ldeen{Liest Rahmen-/Füllfarbe, Dicke und Eckenradius aus der
% \code{enclosure}-Komponente der benannten osgstyler-Dekoration #1;
% \code{pattern} und \code{padding} werden derzeit nicht
% ausgewertet.}{Reads frame/fill color, thickness, and corner radius
% from the named osgstyler decoration #1's \code{enclosure} component;
% \code{pattern} and \code{padding} are not evaluated yet.}
\cs_new_protected:Npn \__semcat_render_box_from_decoration:n #1
  {
    \tl_set:Nx \l__semcat_tmp_tl { \OsgDecorationValue {#1} { enclosure } { border-color } }
    \tl_if_empty:NTF \l__semcat_tmp_tl
      { \tl_set_eq:NN \l__semcat_enclosure_frame_tl \l__semcat_box_color_tl }
      { \tl_set:Nx \l__semcat_enclosure_frame_tl { \OsgColorName { \l__semcat_tmp_tl } } }
    \tl_set:Nx \l__semcat_tmp_tl { \OsgDecorationValue {#1} { enclosure } { fill-color } }
    \tl_if_empty:NTF \l__semcat_tmp_tl
      { \tl_set:Nx \l__semcat_enclosure_back_tl { \l__semcat_box_color_tl!12 } }
      { \tl_set:Nx \l__semcat_enclosure_back_tl { \OsgColorName { \l__semcat_tmp_tl } } }
    \tl_set:Nx \l__semcat_tmp_tl { \OsgDecorationValue {#1} { enclosure } { thickness } }
    \tl_if_empty:NTF \l__semcat_tmp_tl
      { \tl_set:Nn \l__semcat_enclosure_rule_tl { .4pt } }
      { \exp_args:No \__semcat_metric_resolve:nN \l__semcat_tmp_tl \l__semcat_enclosure_rule_tl }
    \tl_set:Nx \l__semcat_tmp_tl { \OsgDecorationValue {#1} { enclosure } { radius } }
    \tl_if_empty:NTF \l__semcat_tmp_tl
      { \tl_set:Nn \l__semcat_enclosure_radius_tl { 3pt } }
      { \exp_args:No \__semcat_metric_resolve:nN \l__semcat_tmp_tl \l__semcat_enclosure_radius_tl }
    \__semcat_compose_title:
    \begin{tcolorbox}[
        enhanced, breakable,
        colback = \l__semcat_enclosure_back_tl,
        colframe = \l__semcat_enclosure_frame_tl,
        fonttitle = \sffamily,
        title = {\l__semcat_box_title_tl},
        arc = \l__semcat_enclosure_radius_tl, boxrule = \l__semcat_enclosure_rule_tl,
        \l__semcat_box_options_tl
      ]
      \BODY
    \end{tcolorbox}
  }

% \ldeen{Entscheidet die tatsächliche Darstellung: Ist eine
% osgstyler-Dekoration aufgelöst, bestimmt deren \code{margin}- bzw.
% \code{enclosure}-Komponente Balken oder Box; ohne Dekoration
% entscheidet der einfache @1-Schlüssel zwischen den eingebauten
% Balken- und Box-Darstellungen.}{Decides the actual rendering: if an
% osgstyler decoration was resolved, its \code{margin} or
% \code{enclosure} component determines a bar or a box; without a
% decoration, the plain @1 key chooses between the built-in bar and box
% renderings.}{\code{style}}
\cs_new_protected:Npn \__semcat_render_mark:
  {
    \tl_if_empty:NTF \l__semcat_resolved_decoration_tl
      {
        \str_if_eq:VnTF \l__semcat_resolved_style_tl { bar }
          { \__semcat_render_bar_default: }
          { \__semcat_render_box_default: }
      }
      {
        \IfPackageLoadedTF { osgstyler }
          {
            \exp_args:No \OsgDecorationIfHasComponentTF \l__semcat_resolved_decoration_tl { margin }
              { \exp_args:No \__semcat_render_bar_from_decoration:n \l__semcat_resolved_decoration_tl }
              {
                \exp_args:No \OsgDecorationIfHasComponentTF \l__semcat_resolved_decoration_tl { enclosure }
                  { \exp_args:No \__semcat_render_box_from_decoration:n \l__semcat_resolved_decoration_tl }
                  { \__semcat_render_box_default: }
              }
          }
          { \__semcat_render_box_default: }
      }
  }

% \ldeen*{Baut auf \cs{NewEnviron} statt auf \cs{NewDocumentEnvironment}
% auf: Der Begin-/End-Code von xparse läuft um den Inhalt zwischen
% \cs{begin}/\cs{end} herum, kann diesen Inhalt aber nicht
% unterdrücken -- eine Moduseinschränkung, die die Markierung
% ausblendet, würde ihren Inhalt trotzdem als nackten Absatz setzen.
% \cs{NewEnviron} fängt den Inhalt zunächst in @1 ab, sodass ein
% inaktiver Modus ihn einfach nie aufruft. Die Umgebung ersetzt sowohl
% das frühere, nur einen Absatz erfassende \cs{SemCatMark} als auch die
% reine Box-Umgebung: welche Darstellung tatsächlich erscheint, legt
% \cs{__semcat_render_mark:} anhand der aufgelösten Kategorie und des
% aktiven Modus fest.}{Built on \cs{NewEnviron} rather than
% \cs{NewDocumentEnvironment}: xparse's begin/end code runs around
% whatever is between \cs{begin}/\cs{end}, but never gets to *suppress*
% that material -- a mode restriction that hides the markup would still
% typeset its body as a bare paragraph. \cs{NewEnviron} captures the
% body into @1 first, so an inactive mode can simply never invoke it.
% The environment replaces both the former, single-paragraph-only
% \cs{SemCatMark} command and the plain box environment: which
% rendering actually appears is decided by \cs{__semcat_render_mark:}
% from the resolved category and the active mode.}{\cs{BODY}}
\bool_new:N \l__semcat_box_shown_bool
\tl_new:N \l__semcat_box_mode_tl
\tl_new:N \l__semcat_box_category_key_tl
\tl_new:N \l__semcat_box_options_tl
\tl_new:N \l__semcat_box_category_tl
\tl_new:N \l__semcat_box_color_tl

\keys_define:nn { semcat / box }
  {
    mode      .tl_set:N  = \l__semcat_box_mode_tl,
    mode      .initial:n = { },
    category  .tl_set:N  = \l__semcat_box_category_key_tl,
    category  .initial:n = { },
    options   .tl_set:N  = \l__semcat_box_options_tl,
    options   .initial:n = { },
  }

\NewEnviron{ SemCatMark }[1][]
  {
    \tl_clear:N \l__semcat_box_mode_tl
    \tl_clear:N \l__semcat_box_category_key_tl
    \tl_clear:N \l__semcat_box_options_tl
    \keys_set:nn { semcat / box } {#1}
    \__semcat_resolve_category:nN \l__semcat_box_category_key_tl \l__semcat_box_category_tl
    \__semcat_validate_category:N \l__semcat_box_category_tl
    \exp_args:No \__semcat_resolve_color:nN \l__semcat_box_category_tl \l__semcat_box_color_tl
    \bool_set_true:N \l__semcat_box_shown_bool
    \tl_if_empty:NF \l__semcat_box_mode_tl
      {
        \cs_if_exist:NT \IfLectureModeTF
          {
            \exp_args:No \IfLectureModeF \l__semcat_box_mode_tl
              { \bool_set_false:N \l__semcat_box_shown_bool }
          }
      }
    \bool_if:NT \l__semcat_box_shown_bool
      {
        \exp_args:No \__semcat_resolve_presentation:nNNN \l__semcat_box_category_tl
          \l__semcat_resolved_style_tl \l__semcat_resolved_decoration_tl
          \l__semcat_resolved_symbol_tl
        \__semcat_resolve_affix:
        \exp_args:No \__semcat_tag_begin_note:n \l__semcat_box_category_tl
        \savenotes
        \__semcat_render_mark:
        \spewnotes
        \__semcat_tag_end:
      }
  }

% ---------------------------------------------------------------------
% QR annotations
% ---------------------------------------------------------------------
\tl_new:N \l__semcat_redirect_base_tl
\tl_new:N \l__semcat_redirect_docid_tl
\tl_set:Nx \l__semcat_redirect_docid_tl { \c_sys_jobname_str }

\NewDocumentCommand \SemCatSetRedirectBase { m O{\c_sys_jobname_str} }
  {
    \tl_set:Nn \l__semcat_redirect_base_tl {#1}
    \tl_set:Nn \l__semcat_redirect_docid_tl {#2}
  }

\cs_new_protected:Npn \__semcat_qr_box:nnn #1#2#3
  { % #1 = url  #2 = caption  #3 = resolved color name
    {
      \qrset{height=1.2cm}
      \cs_if_exist:NT \hypersetup { \hypersetup{urlcolor=#3} }
      \parbox{1.5cm}{%
        \qrcode{#1}\\%
        \tiny\sffamily #2%
      }%
    }%
  }

\cs_new_protected:Npn \__semcat_qr:nnnnnn #1#2#3#4#5#6
  { % #1 mode #2 star("true"/"false") #3 category #4 url #5 caption #6 shift
    \__semcat_if_mode:nn {#1}
      {
        \group_begin:
        \__semcat_resolve_category:nN {#3} \l_tmpa_tl
        \__semcat_validate_category:N \l_tmpa_tl
        \exp_args:No \__semcat_resolve_color:nN \l_tmpa_tl \l__semcat_tmp_color_tl
        % \ldeen{Hier findet keine Struktureinbettung statt: qrcode
        % zeichnet den Code über eine eigene pgf-/tikz-Grafik, und das
        % manuelle Öffnen einer Tag-Struktur um fremdes tikz-Ergebnis
        % erzeugt ungültige Eltern-Kind-Beziehungen im Tag-Baum. Die
        % Beschriftung trägt den eigentlichen semantischen Gehalt und
        % wird als Teil des umgebenden Fließtexts ganz normal
        % getaggt.}{No structure wrapping happens here: qrcode draws
        % the code via its own pgf/tikz picture, and manually opening a
        % tag structure around third-party tikz output produces invalid
        % Parent-Child relations in the tag tree. The caption carries
        % the actual semantic content and is tagged normally as part of
        % the surrounding flow.}
        \str_if_eq:nnTF {#2} { true }
          { \__semcat_qr_box:nnn {#4} {#5} { \l__semcat_tmp_color_tl } }
          {
            \marginnote[{\__semcat_qr_box:nnn {#4} {#5} { \l__semcat_tmp_color_tl }}]
              {\raggedleft\__semcat_qr_box:nnn {#4} {#5} { \l__semcat_tmp_color_tl }}[#6]
          }
        \group_end:
      }
  }

\NewDocumentCommand \SemCatQR { D<>{} s O{} m m O{0pt} }
  {
    \IfBooleanTF {#2}
      { \__semcat_qr:nnnnnn {#1} { true } {#3} {#4} {#5} {#6} }
      { \__semcat_qr:nnnnnn {#1} { false } {#3} {#4} {#5} {#6} }
  }

\cs_new:Npn \__semcat_urlencode:n #1
  {
    \sys_if_engine_luatex:TF
      { \directlua{ tex.print(semcat.urlencode("\luaescapestring{#1}")) } }
      { \msg_warning:nnn { semcat } { luatex-required-for-redirect } {#1} #1 }
  }

\NewDocumentCommand \SemCatQRRedirect { D<>{} s O{} m m O{0pt} }
  {
    \tl_if_empty:NTF \l__semcat_redirect_base_tl
      { \msg_error:nn { semcat } { redirect-base-missing } }
      {
        \IfBooleanTF {#2}
          { \__semcat_qr:nnnnnn {#1} { true } {#3}
              { \l__semcat_redirect_base_tl
                ?doc=\__semcat_urlencode:n{\l__semcat_redirect_docid_tl}
                &id=\__semcat_urlencode:n{#4} }
              {#5} {#6} }
          { \__semcat_qr:nnnnnn {#1} { false } {#3}
              { \l__semcat_redirect_base_tl
                ?doc=\__semcat_urlencode:n{\l__semcat_redirect_docid_tl}
                &id=\__semcat_urlencode:n{#4} }
              {#5} {#6} }
      }
  }

% ---------------------------------------------------------------------
% Floating explanation box: always a box (never a bar), and always a
% genuine LaTeX float (numbered, captioned, page-placement-flexible,
% collectible into a "List of Explanations") via the \pkg{float}
% package's \code{semcatexplain} type -- kept independent from
% \cs{SemCatMark}'s style/decoration resolution, since being a float is
% the defining trait here, not a look that a category or mode could
% swap out.
% ---------------------------------------------------------------------
\NewEnviron{ SemCatExplain }[2][]
  {
    \tl_clear:N \l__semcat_box_mode_tl
    \tl_clear:N \l__semcat_box_category_key_tl
    \tl_clear:N \l__semcat_box_options_tl
    \keys_set:nn { semcat / box } {#1}
    \__semcat_resolve_category:nN \l__semcat_box_category_key_tl \l__semcat_box_category_tl
    \__semcat_validate_category:N \l__semcat_box_category_tl
    \exp_args:No \__semcat_resolve_color:nN \l__semcat_box_category_tl \l__semcat_box_color_tl
    \bool_set_true:N \l__semcat_box_shown_bool
    \tl_if_empty:NF \l__semcat_box_mode_tl
      {
        \cs_if_exist:NT \IfLectureModeTF
          {
            \exp_args:No \IfLectureModeF \l__semcat_box_mode_tl
              { \bool_set_false:N \l__semcat_box_shown_bool }
          }
      }
    \bool_if:NT \l__semcat_box_shown_bool
      {
        \exp_args:No \__semcat_resolve_presentation:nNNN \l__semcat_box_category_tl
          \l__semcat_resolved_style_tl \l__semcat_resolved_decoration_tl
          \l__semcat_resolved_symbol_tl
        \__semcat_resolve_affix:
        \__semcat_tag_begin:n { \UseStructureName{ semcat/explain } }
        \begin{semcatexplain}[tbp]
          \centering
          \begin{tcolorbox}[
              enhanced, breakable,
              colback = \l__semcat_box_color_tl!8,
              colframe = \l__semcat_box_color_tl!70!black,
              fonttitle = \sffamily,
              arc = 3pt, boxrule = .4pt,
              \l__semcat_box_options_tl
            ]
            \BODY
          \end{tcolorbox}
          \tl_if_empty:NTF \l__semcat_affix_content_tl
            { \caption{#2} }
            { \caption{ \textcolor{\l__semcat_affix_color_tl}{\l__semcat_affix_content_tl} ~ #2 } }
        \end{semcatexplain}
        \__semcat_tag_end:
      }
  }

\sys_if_engine_luatex:T { \directlua{ semcat = require("semcat") } }

\ExplSyntaxOff

% ---------------------------------------------------------------------
% Built-in default categories (mirrors the legacy A/B/C/D scheme)
% ---------------------------------------------------------------------
\SemCatDefine*{A}[label={Required}]
\SemCatDefine{B}[label={Supplementary}]
\SemCatDefine{C}[label={Background}]
\SemCatDefine{D}[label={Comparison}]
%</package>
